\documentclass[11pt,a4paper]{article}

% -------------------------------------------------
% Encoding and fonts
% -------------------------------------------------
\usepackage[T1]{fontenc}
\usepackage[utf8]{inputenc}
\usepackage{lmodern}

% -------------------------------------------------
% Page layout
% -------------------------------------------------
\usepackage[
    a4paper,
    top=25mm,
    bottom=25mm,
    left=25mm,
    right=25mm,
    headheight=14pt
]{geometry}

\usepackage{setspace}
\setstretch{1.08}

\usepackage{parskip}
\setlength{\parindent}{0pt}
\setlength{\parskip}{0.65em}

% -------------------------------------------------
% Mathematics
% -------------------------------------------------
\usepackage{amsmath}
\usepackage{amssymb}
\usepackage{mathtools}

% -------------------------------------------------
% Tables
% -------------------------------------------------
\usepackage{array}
\usepackage{booktabs}
\usepackage{tabularx}
\usepackage{longtable}
\usepackage{multirow}
\usepackage{makecell}

% Column definitions for tables
\newcolumntype{L}[1]{>{\raggedright\arraybackslash}p{#1}}
\newcolumntype{C}[1]{>{\centering\arraybackslash}p{#1}}
\newcolumntype{R}[1]{>{\raggedleft\arraybackslash}p{#1}}

% Allow tables to be slightly more compact
\renewcommand{\arraystretch}{1.15}

% -------------------------------------------------
% Quotations and lists
% -------------------------------------------------
\usepackage{csquotes}
\usepackage{enumitem}

\setlist[itemize]{
    leftmargin=2em,
    itemsep=0.2em,
    topsep=0.3em
}

\setlist[enumerate]{
    leftmargin=2em,
    itemsep=0.2em,
    topsep=0.3em
}

% -------------------------------------------------
% Colors and boxes
% -------------------------------------------------
\usepackage{xcolor}

\definecolor{reviewergray}{RGB}{245,245,245}
\definecolor{darkblue}{RGB}{30,70,120}
\definecolor{placeholderred}{RGB}{150,30,30}

% Reviewer quotation environment
\usepackage{framed}

\newenvironment{reviewercomment}
{
    \begin{quote}
    \itshape
    \color{black}
}
{
    \end{quote}
}

% Optional shaded reviewer comment environment
\newenvironment{reviewerbox}
{
    \begin{quote}
    \setlength{\fboxsep}{8pt}
    \colorbox{reviewergray}{%
        \begin{minipage}{0.92\linewidth}
        \itshape
}
{
        \end{minipage}
    }
    \end{quote}
}

% -------------------------------------------------
% Code, paths, and placeholders
% -------------------------------------------------
\usepackage{verbatim}
\usepackage{listings}

\lstset{
    basicstyle=\ttfamily\small,
    breaklines=true,
    breakatwhitespace=false,
    columns=fullflexible,
    keepspaces=true,
    showstringspaces=false
}

% A robust command for inline code and file paths
\newcommand{\code}[1]{\texttt{\detokenize{#1}}}

% Placeholder command
\newcommand{\placeholder}[1]{%
    \textcolor{placeholderred}{\texttt{[PLACEHOLDER: #1]}}%
}

% Optional reviewer tag command
\newcommand{\reviewertag}[1]{%
    \texttt{#1}%
}

% -------------------------------------------------
% Units and scientific notation
% -------------------------------------------------
\usepackage{siunitx}

\sisetup{
    separate-uncertainty=true,
    uncertainty-separator={\,},
    exponent-product=\cdot,
    output-exponent-marker=\mathrm{e},
    range-phrase={--},
    range-units=single
}

% -------------------------------------------------
% Figures and captions
% -------------------------------------------------
\usepackage{graphicx}
\usepackage{caption}
\usepackage{subcaption}

\captionsetup{
    font=small,
    labelfont=bf,
    justification=raggedright,
    singlelinecheck=false
}

% -------------------------------------------------
% References and hyperlinks
% -------------------------------------------------
\usepackage[
    colorlinks=true,
    linkcolor=darkblue,
    citecolor=darkblue,
    urlcolor=darkblue,
    pdftitle={Response to Reviewer 1},
    pdfauthor={}
]{hyperref}

\usepackage[nameinlink,noabbrev]{cleveref}

% -------------------------------------------------
% Typography
% -------------------------------------------------
\usepackage{microtype}
\usepackage{url}

% -------------------------------------------------
% Section formatting
% -------------------------------------------------
\usepackage{titlesec}

\titleformat{\section}
    {\normalfont\Large\bfseries}
    {}
    {0pt}
    {}

\titleformat{\subsection}
    {\normalfont\large\bfseries}
    {}
    {0pt}
    {}

\titleformat{\subsubsection}
    {\normalfont\normalsize\bfseries}
    {}
    {0pt}
    {}

\titlespacing*{\section}
    {0pt}{1.2em}{0.5em}

\titlespacing*{\subsection}
    {0pt}{1.0em}{0.3em}

\titlespacing*{\subsubsection}
    {0pt}{0.8em}{0.2em}

% -------------------------------------------------
% Header and footer
% -------------------------------------------------
\usepackage{fancyhdr}

\pagestyle{fancy}
\fancyhf{}

\lhead{\small Response to Reviewer 1}
\rhead{\small Revised Working Draft}
\cfoot{\thepage}

\renewcommand{\headrulewidth}{0.4pt}
\renewcommand{\footrulewidth}{0pt}

% -------------------------------------------------
% Document metadata
% -------------------------------------------------
\title{\textbf{Response to Reviewer 1}}
\author{}
\date{Draft date: 2026-09-10}

% -------------------------------------------------
% Useful mathematical macros
% -------------------------------------------------
\newcommand{\DFT}{\mathrm{DFT}}
\newcommand{\MLFF}{\mathrm{MLFF}}
\newcommand{\fixed}{\mathrm{fixed}}
\newcommand{\self}{\mathrm{self}}
\newcommand{\RMSE}{\mathrm{RMSE}}
\newcommand{\MAE}{\mathrm{MAE}}
\newcommand{\Rtwo}{R^{2}}
\newcommand{\eV}{\mathrm{eV}}

% -------------------------------------------------
% Begin document
% -------------------------------------------------
\begin{document}

\maketitle
\thispagestyle{fancy}

\section*{Response to Reviewer 1}

\noindent\textbf{Draft date:} 2026-09-10

\vspace{0.5em}

We thank Reviewer~1 for the careful and technically detailed assessment. The comments identified several places where the original manuscript and released code did not sufficiently document the provenance, comparability, or reproducibility of the reported results. We have treated these comments as an opportunity to audit the underlying calculations, clarify the scope of the claims, and strengthen both the manuscript and the code release.

This remains a working response letter. Values already verified in the Revision~1 audit are reported explicitly. Items that require additional cluster calculations, retraining, or final integration remain indicated by bracketed placeholders; their calculation routes are tracked in \texttt{pending\_values\_register.csv}.

\subsection*{A1. Fig.~3d versus Fig.~S3: NEB barrier discrepancy}

\noindent\textbf{Reviewer comment.}
\begin{quote}
\itshape
``Reconcile the foundation-model in-gap barrier discrepancy: state explicitly whether Fig. 3d (fixed-geometry, $\sim$0.7 eV overestimate, `unacceptable') and Fig. S3 (self-run NEB, 0.41 eV, `exceptional') use the same pathway and DFT reference; fix the reference value (0.34 vs.\ $\sim$0.3 eV); present both evaluations side-by-side and explain why relaxation changes the verdict. As written, the paper's central ranking is ambiguous.''
\end{quote}

\noindent\textbf{Response.}
We agree that the original manuscript did not distinguish these two calculations clearly enough. The two figures report different NEB-based evaluation protocols and therefore do not represent the same quantity.

\begin{itemize}
    \item \textbf{Fixed-geometry evaluation on the DFT NEB path (Fig.~3d).}
    The DFT-relaxed NEB images are held fixed, and each MLFF is used only to evaluate the energies of those same geometries. This tests the energetic accuracy of an MLFF along a DFT-defined migration path:
    \begin{equation}
    \Delta E^{\mathrm{fixed}}_M(\Gamma_{\mathrm{DFT}})
    =
    \max_i E_M(X_i^{\mathrm{DFT}})
    -
    E_M(X_0^{\mathrm{DFT}}).
    \end{equation}

    \item \textbf{Self-consistent MLFF NEB optimization (Fig.~S3).}
    The MLFF starts from the endpoint structures and relaxes its own elastic band using its own forces and energy landscape:
    \begin{equation}
    \Gamma_M^\star = \operatorname{NEB}_M(X_0,X_1),
    \qquad
    \Delta E^{\mathrm{self}}_M
    =
    \max_i E_M(X_i^{M,\star})
    -
    E_M(X_0).
    \end{equation}
\end{itemize}

Thus, Fig.~3d evaluates whether a model reproduces the barrier along a fixed DFT path, whereas Fig.~S3 evaluates whether the model can generate a stable migration path under model-in-the-loop NEB optimization. To better clarify this distinction, additional language has been added to explain the purpose and goal of each test.

For the fixed-geometry in-gap pathway 1--4, the archived QE \texttt{neb.out} calculation gives a DFT barrier of 0.336050~eV. However, this archived NEB calculation does not yet meet our final convergence criterion. We therefore treat this value as provisional: it will either be retained with an explicit convergence qualification or replaced by \texttt{[PLACEHOLDER: final-QE-1-4-barrier]} after completion of the restartable QE NEB closure calculation.

Using the current audited pathway~1--4 reference, the fixed-geometry results are:

\begin{table}[h]
\centering
\caption{Historical versus current grouped-split fixed-geometry results for pathway~1--4.}
\begin{tabular}{lcc}
\toprule
Quantity & Value & Interpretation \\
\midrule
DFT reference candidate & 0.336050~eV & provisional \texttt{neb.out} value \\
Historical MACE FT-600K barrier & 0.496792~eV & frame-level, single model \\
Historical absolute error & 0.160742~eV & $0.496792-0.336050$; not reproduced \\
Grouped OMAT-FT absolute error & $0.474\pm0.121$~eV & seeds 123/234/345 \\
Grouped Scratch absolute error & $1.204\pm0.270$~eV & seeds 123/234/345 \\
Grouped Scratch-5\% absolute error & $0.678\pm0.195$~eV & seeds 123/234/345 \\
Grouped MPA0 comparator & 0.237~eV & single seed; no uncertainty estimate \\
\bottomrule
\end{tabular}
\end{table}

The historical $0.16$~eV value is therefore retained only to explain the reviewer comment; it is not substituted for, or identified with, the current DFT reference of $0.336050$~eV. The complete seed-level comparison, including explicit placeholders for the unresolved deep-penetration and grouped FT-Multi$_T$ measurements, is maintained in the grouped comparison table used by the manuscript/SI.

The 0.41~eV value in SI Fig.~S3 is now described as the Foundation model's self-consistent MLFF-NEB barrier rather than a fixed-path DFT-reference error, for better clarity. We have also removed language implying that this result is ``exceptional.'' Instead, the revised text will describe it more narrowly as evidence that the Foundation model can produce a stable path under its own NEB optimization protocol.

As an additional matched fixed-path check, we evaluated ten Foundation checkpoints (MP-0b3, MPA0, OMAT-0 small/medium, MH-0/MH-1, Agnesi medium, and three Agnesi stress/density variants) on the same five pathway-1--4 images. The four representative checkpoints MP-0b3, MPA0, OMAT-0 medium, and MH-1 gave barriers of 1.503, 1.791, 1.602, and 1.841~eV, respectively, relative to the provisional 0.336050~eV DFT candidate. Their absolute errors therefore range from 1.17 to 1.50~eV, and MPA0 is not the kinetic-accuracy winner in this completed comparison. The ten individual profiles and family-aggregated profiles are added to the SI. The family bands represent sub-model variation, not seed uncertainty; all foundation checkpoints are single-checkpoint comparisons. The complete preflight manifest, raw profiles, profile diagnostics, and SHA256 ledger are released with the revision artifacts.

We have added a side-by-side SI table reporting, for each NEB result, the pathway, endpoint definition, image source, relaxation engine, energy engine, barrier definition, DFT reference where applicable, and convergence status.

\noindent\textbf{Manuscript/SI change.}
Main-text Sec.~3.2 and the SI Fig.~S3 text and caption are marked with reviewer tags \texttt{REV-A1} and \texttt{REV-C1}. The proposed replacement text is provided in \texttt{reviewer\_a1\_fig3\_figs3\_update/proposed\_text\_edits\_E3\_E4\_SI\_replacement.md}.

\subsection*{A2. Train/test overlap for FT-600K}

\noindent\textbf{Reviewer comment.}
\begin{quote}
\itshape
``Rule out train/test overlap for FT-600K: report the minimum SOAP-distance (or RMSD) between each in-gap NEB image and its nearest fine-tuning frame, or re-run FT-600K with the NEB-adjacent trajectory segment excluded, and confirm the 0.16 eV result survives.''
\end{quote}

\noindent\textbf{Response.}
We agree that this check is necessary. We have revised the data-splitting procedure so that structures are partitioned by trajectory or pathway provenance rather than by an independent frame-level shuffle: every source trajectory is assigned in full to train, validation, or test (1,786/311/267 frames), so no trajectory is split across partitions.

The completed quantitative audit is:

\begin{center}
Full grouped training set (1,786 frames): minimum SOAP cosine distance to any in-gap NEB image $= 5.4\times10^{-4}$; minimum comparable Cartesian RMSD $= 2.06$~\AA. Scratch-5\% subset (89 frames): minimum SOAP cosine distance $= 6.4\times10^{-4}$; minimum comparable RMSD $= 2.06$~\AA. Both are far outside the exact/near-duplicate thresholds used for this audit (SOAP $\le10^{-4}$, RMSD $\le0.05$~\AA); zero exact or near-duplicate matches were found for either split.
\end{center}

We additionally retrained FT-600K (from the OMAT-FT and MACE-MPA-0 checkpoints), Scratch, and Scratch-5\% under this grouped split across seeds 123/234/345 (MPA0: single seed, excluded from seed-uncertainty statistics). On the in-gap pathway 1--4, the absolute fixed-path barrier error relative to the current DFT candidate (0.336050~eV) is $0.474 \pm 0.121$~eV for FT-600K-OMATFT, $1.204 \pm 0.270$~eV for Scratch, $0.678 \pm 0.195$~eV for Scratch-5\%, and $0.237$~eV for the single-seed MPA0 comparator. FT-600K-OMATFT remains the smallest-error fine-tuned model on this pathway, consistent with the original single-model result, but the historical 0.16~eV point estimate is not reproduced under this grouped, multi-seed re-evaluation. We report this directly rather than asserting that the 0.16~eV result survives. The remaining four audited pathways (1--2, 1--3, 1--5, 1--7) were also evaluated under the grouped protocol but are withheld from this comparison pending confirmation of their benchmark classification and DFT-reference identity; retraining FT-Multi\_T under the grouped protocol was not part of this round.

\noindent\textbf{Manuscript change.}
Methods Sec.~2.1 and the in-gap diffusion results (Sec.~3.3.3) are marked with \texttt{REV-A2/REV-C10}, and Table~S1 in the SI reports the grouped-split RMSE and barrier statistics.

\subsection*{A3. Uncertainty in kinetic results}

\noindent\textbf{Reviewer comment.}
\begin{quote}
\itshape
``Add uncertainty to the kinetic results: retrain at least Scratch and FT-600K with $\geq$3 seeds and report barrier-error mean +/- std (the 3-seed protocol already exists for Table S1). This simultaneously tests the `chance' explanation for the scratch model's deep-penetration result---if the low error is seed-stable, `chance' is refuted; if seed-dependent, it is supported. Either way, replace the current unfalsifiable assertion with evidence.''
\end{quote}

\noindent\textbf{Response.}
We agree. We have converted this issue into a targeted robustness assessment. Scratch and FT-600K were trained using seeds 123, 234, and 345 under the revised grouped-split protocol; this training is complete.

For the \emph{in-gap} pathway (1--4), the absolute barrier errors are $0.474 \pm 0.121$~eV (FT-600K-OMATFT) and $1.204 \pm 0.270$~eV (Scratch) across the three seeds (individual seed values: FT-600K-OMATFT 0.613/0.413/0.396~eV; Scratch 1.516/1.036/1.061~eV). For the \emph{deep-penetration} pathway, the multi-seed models have been trained and evaluated on all five audited pathways, but we have not yet confirmed which of the remaining pathway IDs corresponds to the deep-penetration probe used in the original submission, nor closed out the restarted QE reference for it:

\begin{center}
\texttt{[PLACEHOLDER: deep-penetration pathway classification, converged DFT reference, and seed-level barrier errors]}
\end{center}

The original ``by chance'' wording has been removed. For the in-gap pathway, the multi-seed evidence does not support an unqualified ``by chance'' dismissal or an unqualified robustness claim: FT-600K-OMATFT is the smallest-error model with a seed spread of $\pm0.121$~eV, and Scratch shows a comparable seed spread ($\pm0.270$~eV) at a higher absolute error. For the deep-penetration pathway specifically, the revised manuscript will state only what is supported once the pathway classification and DFT reference above are resolved.

\noindent\textbf{Manuscript change.}
The Sec.~3.2 deep-penetration discussion and Table~S1 are marked with \texttt{REV-A3}.

\subsection*{A4. Scope of generalizability claims}

\noindent\textbf{Reviewer comment.}
\begin{quote}
\itshape
``Rescope the generalizability claims: change `establishes \ldots{} a generalizable standard' to a candidate framework demonstrated on one system, and move `broader principles' to a hedged outlook---or add one additional material/dopant (even a cheap second system, e.g., a different transition-metal dopant in the same host) as minimal cross-validation.''
\end{quote}

\noindent\textbf{Response.}
We agree. The revised manuscript no longer presents the workflow as a general standard. Instead, it describes the study as a candidate evaluation framework demonstrated for Cr-doped Sb$_2$Te$_3$ using the MACE architecture family. Broader generality is framed as an outlook that requires validation across additional chemistries, defect environments, and MLIP architectures.

No additional numerical result is required for this revision unless we decide to include an optional second-dopant validation case:

\begin{center}
\texttt{[OPTIONAL PLACEHOLDER: second-dopant-cross-validation]}
\end{center}

\noindent\textbf{Manuscript change.}
The abstract, introduction, research questions, and conclusion are marked with \texttt{REV-A4}.

\subsection*{A5. Scratch-5\% control}

\noindent\textbf{Reviewer comment.}
\begin{quote}
\itshape
``Add a Scratch-5\% control (same $\sim$1,000 configurations as FT-600K, random initialization) to disentangle foundation-pretraining benefits from data-volume effects.''
\end{quote}

\noindent\textbf{Response.}
We agree that this is the appropriate control. The Scratch-5\% models use the same grouped FT-600K training subset (89 frames), identical MACE hyperparameters, random initialization, and seeds 123, 234, and 345; this training is complete and has been evaluated using equilibrium RMSE and the in-gap migration-barrier task used for the fine-tuned models.

\begin{center}
Test RMSE: $16.7 \pm 9.2$~meV/atom (energy), $341.5 \pm 15.0$~meV/\AA\ (force). Pathway 1--4 absolute barrier error: $0.678 \pm 0.195$~eV, versus $0.474 \pm 0.121$~eV for FT-600K-OMATFT and $1.204 \pm 0.270$~eV for full Scratch.
\end{center}

On this single pathway and with only $\sim$90 training frames, Scratch-5\% falls between FT-600K and full Scratch, which is consistent with, but does not by itself establish, a foundation-pretraining benefit beyond data volume; we report the comparison at this hedged level given the small sample (one pathway, three seeds). The same comparison against FT-MultiT and against the remaining four audited pathways is not yet available because FT-MultiT was not retrained under the grouped protocol and those pathways await benchmark-classification confirmation.

\noindent\textbf{Manuscript change.}
Sec.~3.3.3 and Table~S1 are marked with \texttt{REV-A5}.

\subsection*{A6. Interpretability claims: t-SNE/PHATE and SHAP}

\noindent\textbf{Reviewer comment.}
\begin{quote}
\itshape
``Soften the interpretability claims: reframe t-SNE/PHATE conclusions as `consistent with' rather than `direct mechanistic explanation,' report projection sensitivity to perplexity/seed, and compute at least one separation metric in the original descriptor space. For SHAP, report per-model surrogate $R^2$, rescope `the model internalizes X' to `the surrogate's error predictions are most sensitive to X,' and (ideally) add a perturbation test on MACE's own inputs for the dominant Cr-Cr channel.''
\end{quote}

\noindent\textbf{Response.}
We agree. We no longer describe t-SNE or PHATE projections as direct mechanistic explanations. The revised text presents these analyses as descriptive visualizations that are consistent with, but do not independently establish, the observed model behavior. We now report quantitative separation metrics in the original descriptor space and explicitly document their dependence on feature scaling and projection choices.

The local re-analysis gives:

\begin{table}[h]
\centering
\caption{Silhouette-score analysis across projected and descriptor spaces.}
\begin{tabular}{p{0.50\linewidth}rp{0.27\linewidth}}
\toprule
Space & Overall silhouette & Notes \\
\midrule
Old 2D t-SNE, $p=30$, seed 42 & 0.409 & Retained only as the original projection-space baseline \\
2D t-SNE, $p \in \{5,30,50\}$, 3 seeds & $0.397 \pm 0.027$ & Projection-sensitivity analysis \\
PHATE 2D, 3 seeds & $0.473 \pm 0.000$ & Projection-sensitivity analysis \\
Original 6191-d, Euclidean & 0.333 & Primary descriptor-space metric \\
Original 6191-d, cosine & 0.277 & Alternative descriptor-space metric \\
Original 6191-d, z-scored Euclidean & -0.019 & Demonstrates scaling sensitivity \\
PCA 95\%, z-scored Euclidean & -0.026 & Conservative high-dimensional check \\
\bottomrule
\end{tabular}
\end{table}

For SHAP, we now distinguish in-sample fit from cross-validated surrogate fidelity. The revised 5-fold cross-validation metrics are:

\begin{table}[h]
\centering
\caption{Cross-validated surrogate-model performance used for SHAP analysis.}
\begin{tabular}{lrr}
\toprule
Model & 5-fold CV $R^2$ & 5-fold CV MAE \\
\midrule
FT-600K & 0.9819 & 0.0380 \\
FT-MultiT & 0.8756 & 0.0269 \\
Scratch & 0.9730 & 0.0262 \\
\bottomrule
\end{tabular}
\end{table}

These values support the use of the surrogate as a model of the selected error target. However, the revised text will state explicitly that SHAP explains the surrogate model's predictions of error, rather than directly attributing causality within MACE itself. A direct perturbation analysis on MACE inputs remains pending:

\begin{center}
\texttt{[PLACEHOLDER: MACE-CrCr-perturbation-test]}
\end{center}

\noindent\textbf{Manuscript change.}
Sec.~3.4, the Fig.~5 caption, Sec.~3.5, and the Fig.~6 caption are marked with \texttt{REV-A6/REV-C11/REV-C12/REV-C13}.

\subsection*{A7. Related work and distinction from our earlier report}

\noindent\textbf{Reviewer comment.}
\begin{quote}
\itshape
``Position against the nearest neighbors: cite and discuss arXiv:2510.05020 (MACE fine-tuning for Li-ion diffusion), the foundational-MLIP migration-barrier evaluation (RSC Digital Discovery d5dd00534e), and arXiv:2405.07105 (systematic softening---currently invoked as if novel). Add one paragraph explicitly stating this manuscript's delta over the authors' own arXiv:2509.00090 [46].''
\end{quote}

\noindent\textbf{Response.}
We agree. We have added the requested citations and a paragraph explicitly describing the relationship to the earlier report. The revised introduction positions the present work relative to recent studies of MACE fine-tuning for Li-ion diffusion, migration-barrier evaluation using foundation MLIPs, and systematic softening of universal MLIP potential-energy surfaces.

The revised manuscript also clarifies the additions relative to the earlier report, including expanded transport analysis, self-consistent NEB stability analysis, surrogate-based SHAP analysis, and interlayer-sliding analysis. Statements relying on the corrected MD protocol and final reproducibility materials will be updated after those revision items are completed.

\noindent\textbf{Manuscript change.}
The introduction, Sec.~3.2, softening discussion, and bibliography are marked with \texttt{REV-A7}.

\subsection*{C8. Fabricated all-zero forces in NEB exporters}

\noindent\textbf{Reviewer comment.}
\begin{quote}
\itshape
``Fake zero forces in NEB$\rightarrow$training-data exporters (\texttt{examples/Sb2Te3\_Cr\_doped/03\_neb/neb\_data\_collector.py:442-443}, \texttt{neb\_to\_extended\_xyz.py:278-282}): real QE energies paired with fabricated all-zero forces written into training-format extxyz under \texttt{forces:R:3}; identical bug in two scripts. If any such frames fed fine-tuning (\texttt{forces\_weight=100} per repo README), the model was silently trained toward zero forces. Not proven wired into FT-600K/FT-MultiT (merge config references nonexistent folders). $\rightarrow$ Export real forces from QE output or hard-disable; authors must confirm/deny use in reported runs.''
\end{quote}

\noindent\textbf{Response.}
We agree that this was a serious code issue. We restored the two affected exporters and made them fail closed: the collector now terminates when real QE forces are unavailable, while the extended-XYZ exporter defaults to omitting forces and raises an explicit error if force export is requested without validated force data. No fabricated all-zero force array is written by either script.

We then performed a hash-bound, frame-level provenance audit of the datasets used for the reported fine-tuned models. The audit covered 9,163 frames across the historical FT-600K, historical FT-Multi$_T$, grouped FT-600K, and grouped Scratch-5\% train/validation/test exports. Every frame contains a \texttt{forces:R:3} field. Exactly nine frames have forces that are identically zero; all nine are explicitly labelled one-atom \texttt{IsolatedAtom} Cr/Sb/Te reference records used for energy-offset ($E_0$) bookkeeping. There are zero unclassified all-zero-force frames and zero NEB-derived all-zero-force frames. Thus, the audit found no evidence that the fabricated NEB exports were included in the reported training datasets.

The release contains the complete CSV/SQLite manifest and machine-readable summary, including source-file checksums, export-time metadata, split assignments, and the acceptance result. The corresponding code and audit artifacts are in release commit \texttt{f5c6d19} (full commit: \texttt{f5c6d195a448cc1c9d5b7f5177757dc937d28276}).

\noindent\textbf{Manuscript/code change.}
Methods data provenance and Code Availability are marked with \texttt{REV-C8}. No new figure is needed: the compact manifest summary and reproducibility files provide the direct evidence for this code-provenance issue.

\subsection*{C9. Asymmetric MD protocol for naive fine-tuning}

\noindent\textbf{Reviewer comment.}
\begin{quote}
\itshape
``Asymmetric MD protocol for the naive-FT condition (\texttt{results/MD/2-naive-fine-tuning/mace\_md\_eval\_0814.py:359-360}: \texttt{steps=2000}, \texttt{interval=5} vs.\ \texttt{steps=100000}, \texttt{interval=200} in the other three byte-identical scripts): 50$\times$ shorter trajectory with 40$\times$ denser sampling in the cross-model transport comparison; post-processing assumes 200 fs/frame, so $D/\kappa$ for that condition may be off $\sim$40$\times$.''
\end{quote}

\noindent\textbf{Response.}
We agree. The original FT-600K transport value is not retained as a final model comparison because it was obtained under a non-comparable protocol. It has therefore been quarantined pending completion of a corrected matched-protocol calculation.

The current local transport re-analysis is summarized below:

\begin{table}[h]
\centering
\caption{Current status of transport estimates.}
\begin{tabular}{p{0.16\linewidth}p{0.22\linewidth}p{0.22\linewidth}p{0.30\linewidth}}
\toprule
Model & $D$ (cm$^2$\,s$^{-1}$) & $\kappa$ (W\,m$^{-1}$\,K$^{-1}$) & Interpretation \\
\midrule
Foundation &
$(0.7 \pm 7.9)\times10^{-9}$ &
$(2.4 \pm 8.6)\times10^{-3}$ &
Statistically indistinguishable from zero within current uncertainty \\
Scratch &
$(-3.4 \pm 3.2)\times10^{-9}$ &
$(0.9 \pm 1.3)\times10^{-2}$ &
Statistically indistinguishable from zero; the negative fitted estimate reflects finite-sampling uncertainty rather than a physical negative diffusivity \\
FT-MultiT &
$(4.8 \pm 4.8)\times10^{-8}$ &
$(2.0 \pm 1.7)\times10^{-2}$ &
Preliminary estimate under the matched local re-analysis \\
FT-600K &
$(3.9 \pm 4.0)\times10^{-8}$ &
$(0.9 \pm 1.2)\times10^{-2}$ &
Previous protocol non-comparable; quarantined \\
\bottomrule
\end{tabular}
\end{table}

The corrected FT-600K calculation remains pending:

\begin{center}
\texttt{[PLACEHOLDER: corrected-FT600K-MD-transport]}
\end{center}

\noindent\textbf{Manuscript/code change.}
The MD protocol description, transport figure, and SI transport table are marked with \texttt{REV-C9}.

\subsection*{C10. Ungrouped train/test split}

\noindent\textbf{Reviewer comment.}
\begin{quote}
\itshape
``Ungrouped random train/test split (\texttt{split\_data.py:52-66}): frame-level shuffle over merged trajectories with no grouping by pathway/trajectory ID; near-duplicate adjacent frames can land in both train and test, inflating Table S1 test RMSE. Provenance metadata embedded by the collector (\texttt{path\_name}, \texttt{neb\_image}) is never read downstream---no NEB-exclusion safeguard exists.''
\end{quote}

\noindent\textbf{Response.}
We agree. The split has been revised to group structures by trajectory provenance, with explicit zero-overlap checks across training, validation, and test partitions, preventing information leakage through temporally or structurally adjacent configurations.

The grouped-split retraining and evaluation are complete:

\begin{center}
Test RMSE (energy / force, meV/atom / meV/\AA), mean $\pm$ std over seeds 123/234/345: FT-600K-OMATFT $40.7\pm5.0$ / $259.7\pm5.8$; Scratch $30.5\pm3.6$ / $343.6\pm35.9$; Scratch-5\% $16.7\pm9.2$ / $341.5\pm15.0$; MACE-MPA-0 comparator (single seed) $38.7$ / $278.5$.
\end{center}

These grouped-split RMSE values are reported alongside the historical frame-level split values in Table~S1, which is now labelled to distinguish the two splits; quantitative generalization claims in Sec.~3.3.3 use the grouped-split values.

\noindent\textbf{Manuscript/code change.}
The Methods split description and Table~S1 are marked with \texttt{REV-C10}.

\subsection*{C11. Silhouette scores computed on embeddings}

\noindent\textbf{Reviewer comment.}
\begin{quote}
\itshape
``Silhouette scores computed on the 2D t-SNE embedding, not the original feature space (\texttt{results/t-SNE/3d/tsne.py:323}, \texttt{phate-tsne/tsne.py:191}, \texttt{mace\_probe\_blackbox0818.py:254}): the Fig. 5d separability numbers measure projection artifacts, not true latent separability. Need to compute in original/high-dimensional space.''
\end{quote}

\noindent\textbf{Response.}
We agree. The revised analysis reports descriptor-space silhouette values and identifies the previous embedding-space values only as projection-space baselines. The principal descriptor-space result is an Euclidean silhouette score of 0.333 in the original 6191-dimensional feature space. The SI will additionally report cosine, z-scored Euclidean, and PCA-reduced analyses to make the dependence on feature representation explicit.

No numerical placeholder remains for the local analysis, although the manuscript integration remains to be completed in Overleaf.

\noindent\textbf{Manuscript change.}
Sec.~3.4 and the Fig.~5 caption are marked with \texttt{REV-C11}.

\subsection*{C12. Always-zero force-sensitivity feature stub}

\noindent\textbf{Reviewer comment.}
\begin{quote}
\itshape
``Always-zero `force sensitivity' feature stub in all three latent-probe scripts (\texttt{3d/tsne.py:209-215}, \texttt{phate-tsne/tsne.py:132}, \texttt{mace\_probe\_blackbox0818.py:171-180}), silently concatenated into the feature matrix feeding PCA$\rightarrow$t-SNE$\rightarrow$PHATE.''
\end{quote}

\noindent\textbf{Response.}
We confirm this issue. The inactive feature channel has been removed from the revised latent-space analysis. We do not make any force-sensitivity claim based on this channel, and we will not introduce such a claim unless it is supported by a genuine finite-difference or otherwise physically defined force-sensitivity calculation.

No numerical placeholder remains.

\noindent\textbf{Manuscript/code change.}
The SI latent-analysis methods and code-release notes are marked with \texttt{REV-C12}.

\subsection*{C13. SHAP code and repository stubs}

\noindent\textbf{Reviewer comment.}
\begin{quote}
\itshape
``SHAP analysis code absent from the repository entirely; advertised \texttt{migrationbench} package (36 files), \texttt{scripts/} (13 files including \texttt{shap\_analysis.py}), templates, workflow docs, and \texttt{paper/manuscript.tex} are all 0-byte stubs. Section 3.5 results and the `benchmark framework' API are not reproducible from this release.''
\end{quote}

\noindent\textbf{Response.}
We agree. The release now includes the X-FORCE three-model, 150-frame inputs, an explicit image-aligned SOAP/structural feature workflow, the TreeSHAP analysis scripts, pinned dependencies, and the empty-stub audit. The source lineage is recorded in \texttt{docs/audit/shap\_source\_lineage.md}; the four-model, 2-concentration, and Dual-X ID/OOD/NEB workflows are explicitly excluded from the Sec.~3.5 claim.

The locked revised CSVs report the audited cross-validated values used in Sec.~3.5. An independent local rerun reproduces the expected 150-frame/450-row alignment but is not numerically identical to those locked CSVs; this discrepancy is recorded in the lineage audit and the paper outputs are not silently overwritten pending recovery of the exact historical regeneration implementation.

The remaining release-level evidence is:

\begin{center}
\texttt{C13 GitHub release commit: a3a73dfe88aa96416a3bc365909fa8165d42a271}
\end{center}

\begin{center}
\texttt{docs/audit/c13\_empty\_stub\_audit.md; docs/audit/shap\_source\_lineage.md}
\end{center}

\noindent\textbf{Manuscript/code change.}
Sec.~3.5 and Code Availability are marked with \texttt{REV-C13}.

% \subsection*{Status summary}

% The response structure and the principal conceptual corrections are in place. Completed items include the clarification of the Fig.~3d/Fig.~S3 protocol distinction, the scope revisions, the related-work positioning, the descriptor-space and projection-sensitivity re-analysis, and the SHAP surrogate cross-validation audit.

% The main remaining items are final QE NEB convergence, SOAP/RMSD overlap auditing, grouped-split multi-seed retraining, the Scratch-5\% control, corrected matched-protocol FT-600K MD, and final code-release evidence for the exporter and SHAP-pipeline fixes.


\end{document}
